\section{Introduction: the problem, and how we solve it}
\label{sec:intro}

\subsection{The problem}

A \emph{point process} records the times $t_1<t_2<\cdots$ of events. The
\emph{Hawkes process} is the canonical model when events beget events --- an
earthquake makes aftershocks more likely, a tweet makes re-tweets more likely.
Its defining feature is a \emph{self-exciting} conditional intensity
\begin{equation}
\lambda(t)=s(t)+\sum_{t_i<t}\phi(t-t_i),
\label{eq:intro-hawkes}
\end{equation}
where $s\ge 0$ is the exogenous ``background'' rate that injects \emph{immigrant}
events and the kernel $\phi\ge 0$ describes how each past event lifts the future
rate, producing \emph{offspring}. To fit \eqref{eq:intro-hawkes} by maximum
likelihood one needs every event time $t_i$.

In practice we frequently \emph{do not} observe event times. We observe
\emph{interval-censored} data: the \emph{number} of events in each time bucket ---
hospital admissions per day, video views per day, transactions per hour --- but
not when each event occurred, whether for privacy or because that is simply how
the data are recorded. Two facts then block a direct fit:
\begin{enumerate}[leftmargin=2.2em]
\item the log-likelihood of \eqref{eq:intro-hawkes} is \emph{undefined} without
the times $t_i$; and
\item the Hawkes process lacks the \emph{independent-increments} property, so we
cannot fall back on a Poisson likelihood for the bucket counts either (its
intensity, hence the number of events in a bucket, depends on the counts in
earlier buckets).
\end{enumerate}

\subsection{How we solve it}

The escape is to fit a \emph{different but parameter-matched} process whose mean
behaviour mimics the Hawkes process while enjoying independent increments. This is
the \emph{Mean Behavior Poisson process} (MBPP): the inhomogeneous Poisson process
whose deterministic intensity equals the \emph{expected} Hawkes intensity,
\begin{equation}
\xi(t):=\E[\lambda(t)]=s(t)+\int_0^t \phi(t-u)\,\xi(u)\dd u .
\label{eq:intro-mbpp}
\end{equation}
Three observations make this work, and they organise the whole document.
\begin{itemize}[leftmargin=1.4em]
\item \textbf{Why the MBPP is Poisson.} A simple point process whose compensator
is \emph{deterministic} is automatically Poisson (Watanabe's theorem,
\S\ref{sec:poisson-hawkes}); replacing the random Hawkes intensity $\lambda$ by its
mean $\xi$ removes the feedback and manufactures independent increments. The bucket
counts of the MBPP are therefore independent Poisson variables with means
$\Xi(o_{i-1},o_i]=\int_{o_{i-1}}^{o_i}\xi$, and their negative log-likelihood ---
the \emph{interval-censored log-likelihood} --- is a tractable sum of Poisson
terms.
\item \textbf{Why we can compute $\xi$.} Equation~\eqref{eq:intro-mbpp} is a
\emph{Volterra integral equation of the second kind}. Because $\phi$ here depends
only on the elapsed time $t-u$, it is a \emph{convolution}, so the system is
linear and time-invariant and we can solve it in closed form (an impulse
response), in the Laplace/Fourier domain (a transfer function), or as a linear
ODE. Appendix~\ref{app:volterra} is a self-contained primer on Volterra equations;
Appendices~\ref{app:transforms} and \ref{app:ode} cover the transform and ODE
methods.
\item \textbf{Why the parameters transfer.} The MBPP and the Hawkes process are
built from the \emph{same} $s$ and $\phi$, so fitting the MBPP on the counts
recovers (approximations of) the Hawkes parameters.
\end{itemize}

\subsection{Roadmap, in increasing order of difficulty}

We deliberately build up the model one layer at a time.
\begin{description}[leftmargin=1.4em,style=nextline]
\item[\S\ref{sec:poisson-hawkes} Poisson and Hawkes.] The two ingredients:
the Poisson process (and exactly why a deterministic compensator makes a process
Poisson) and the Hawkes process (and exactly why it is \emph{not} Poisson), with
the cluster/branching picture that yields \eqref{eq:intro-mbpp} rigorously.
\item[\S\ref{sec:assumptions} Standing assumptions.] The precise regularity
conditions under which everything below holds, stated once.
\item[\S\ref{sec:no-covariates} Without covariates.] The univariate and
multivariate MBPP. Here $\phi$ is a fixed convolution kernel, the equation is LTI,
and everything is closed-form. This is the workhorse.
\item[\S\ref{sec:baseline-cov} Baseline covariates.] Let the immigrant rate
depend on covariates, $s(t)=\exp(\gamma_0+\gamma^\top X(t))$. The equation stays a
convolution (only the forcing changes), so the LTI machinery survives unchanged;
we simply fit the coefficients $\gamma$.
\item[\S\ref{sec:excitation-cov} Excitation covariates (the hard case).] Let the
\emph{excitation itself} depend on covariates, $\alpha(t)=\alpha_0\exp(\delta^\top
Z(t))$. Now the kernel depends on $t$ and $u$ \emph{separately}, the equation is a
\emph{general (non-convolution) Volterra equation}, and time-invariance ---
together with the Laplace/Fourier shortcuts --- is lost. We show the problem is
still well-posed and, for exponential kernels, reduces to a \emph{linear
time-varying} ODE that we can integrate.
\end{description}

\noindent
The short answer to ``can we put covariates in the excitation?'' is
\emph{yes} --- \S\ref{sec:excitation-cov} does it --- but it is genuinely harder,
and seeing exactly \emph{why} is the reason the Volterra viewpoint is worth
developing carefully.
